\documentclass[9pt,a4paper]{article}

\usepackage[a4paper,top=11mm,bottom=14mm,left=13mm,right=13mm]{geometry}
\usepackage{fontspec}
\usepackage[table]{xcolor}
\usepackage{tabularx}
\usepackage{array}
\usepackage{enumitem}
\usepackage[most]{tcolorbox}
\usepackage{microtype}
\usepackage{fancyhdr}

\IfFontExistsTF{Arial}{
  \setmainfont{Arial}
  \setsansfont{Arial}
}{
  \setmainfont{TeX Gyre Heros}
  \setsansfont{TeX Gyre Heros}
}

\definecolor{AIEPRed}{HTML}{D52B1E}
\definecolor{AIEPNavy}{HTML}{0B2038}
\definecolor{AIEPInk}{HTML}{16293D}
\definecolor{AIEPMuted}{HTML}{526577}
\definecolor{AIEPLine}{HTML}{C9D3DC}
\definecolor{AIEPSoft}{HTML}{F5F1E9}
\definecolor{GeoCyan}{HTML}{007C91}
\definecolor{GeoGreen}{HTML}{2E7D32}
\definecolor{GeoGreenSoft}{HTML}{E9F2EA}
\definecolor{GeoCyanSoft}{HTML}{E7F3F5}
\definecolor{GeoGold}{HTML}{B77900}
\definecolor{GeoGoldSoft}{HTML}{FFF7DF}

\setlength{\parindent}{0pt}
\setlength{\parskip}{3pt}
\setlist[itemize]{leftmargin=4mm,topsep=1.5pt,itemsep=1.4pt,parsep=0pt}
\renewcommand{\arraystretch}{1.23}
\pagestyle{fancy}
\fancyhf{}
\renewcommand{\headrulewidth}{0pt}
\setlength{\footskip}{7mm}
\fancyfoot[L]{\scriptsize\color{AIEPMuted}GeoGreen Escolar · AIEP Osorno × Instituto Comercial Liceo Bicentenario}
\fancyfoot[R]{\scriptsize\color{AIEPMuted}10 de agosto de 2026}

\newcolumntype{Y}{>{\raggedright\arraybackslash}X}
\newcolumntype{C}[1]{>{\centering\arraybackslash}p{#1}}

\newtcolorbox{cleanbox}[2][]{
  enhanced,
  colback=white,
  colframe=AIEPLine,
  boxrule=0.45pt,
  arc=1mm,
  left=3mm,right=3mm,top=2.3mm,bottom=2.3mm,
  title={\scriptsize\bfseries\color{#2}},
  coltitle=#2,
  fonttitle=\sffamily,
  attach title to upper={\par\vspace{1mm}},
  #1
}

\begin{document}

% Firma visual AIEP × GeoGreen
\noindent
\colorbox{AIEPRed}{\rule{15mm}{0pt}\rule{0pt}{2.2mm}}%
\colorbox{GeoCyan}{\rule{15mm}{0pt}\rule{0pt}{2.2mm}}%
\colorbox{GeoGreen}{\rule{15mm}{0pt}\rule{0pt}{2.2mm}}%
\colorbox{AIEPNavy}{\makebox[\dimexpr\textwidth-45mm-8\fboxsep\relax][r]{\scriptsize\bfseries\color{white}COORDINACIÓN 2026\hspace{1mm}}}

\vspace{4mm}
{\scriptsize\bfseries\color{AIEPRed}AIEP OSORNO × GEOGREEN ESCOLAR\par}
\vspace{1.5mm}
{\fontsize{24}{27}\selectfont\bfseries\color{AIEPNavy}Ajuste de programación\par}
\vspace{1mm}
{\large\color{AIEPMuted}Especialidad de Programación · 3° A y 3° C\par}

\vspace{4mm}
\begin{tcolorbox}[
  enhanced,
  colback=GeoGreenSoft,
  colframe=GeoGreen,
  boxrule=0.7pt,
  borderline west={2.2mm}{0pt}{GeoGreen},
  arc=1mm,
  left=4mm,right=4mm,top=2.5mm,bottom=2.5mm]
{\scriptsize\bfseries\color{GeoGreen}RESULTADO DEL AJUSTE}\par
{\normalsize\bfseries\color{AIEPNavy}Se conservan íntegramente los 16 bloques definidos por el establecimiento.}\par
{\small\color{AIEPInk}Mismas fechas, horarios, docentes y duración de 90 minutos; se corrigen solamente errores de registro y secuencia.}
\end{tcolorbox}

\vspace{1mm}
\begin{minipage}[t]{0.487\textwidth}
\begin{cleanbox}[height=38mm]{AIEPNavy}
{\scriptsize\bfseries\color{AIEPNavy}SE CONSERVA}
\begin{itemize}
  \item Los días y bloques horarios propuestos.
  \item La ejecución separada para 3° A y 3° C.
  \item Los docentes asignados a cada bloque.
  \item La duración de 90 minutos por sesión.
\end{itemize}
\end{cleanbox}
\end{minipage}\hfill
\begin{minipage}[t]{0.487\textwidth}
\begin{cleanbox}[height=38mm]{GeoCyan}
{\scriptsize\bfseries\color{GeoCyan}SE CORRIGE}
\begin{itemize}
  \item Todas las fechas quedan registradas en 2026.
  \item El Taller 1 queda el miércoles 19/08.
  \item En 3° C, Taller 3 queda antes de Taller 4.
  \item Se normaliza el formato para una lectura directa.
\end{itemize}
\end{cleanbox}
\end{minipage}

\vspace{2mm}
\begin{tcolorbox}[
  colback=GeoCyanSoft,
  colframe=GeoCyan,
  boxrule=0.55pt,
  arc=1mm,
  left=3mm,right=3mm,top=2mm,bottom=2mm]
\begin{tabularx}{\textwidth}{@{}C{8mm}Y C{10mm} Y@{}}
{\large\bfseries\color{GeoCyan}3} &
{\bfseries\color{AIEPNavy}Taller 3 · Sensores y GeoGreen}\newline
{\small\color{AIEPMuted}Lunes 24 de agosto · ambos cursos}
& {\Large\bfseries\color{GeoCyan}$\longrightarrow$} &
{\bfseries\color{AIEPNavy}Taller 4 · Del dato a la acción}\newline
{\small\color{AIEPMuted}Miércoles 26 de agosto · ambos cursos}
\end{tabularx}
\end{tcolorbox}

\vspace{2mm}
{\scriptsize\bfseries\color{AIEPNavy}MENTORÍAS · MISMA RUTA, CALENDARIOS POR CURSO\par}
\vspace{1.5mm}
\begin{tcolorbox}[
  colback=white,
  colframe=AIEPLine,
  boxrule=0.45pt,
  arc=1mm,
  left=2mm,right=2mm,top=1.5mm,bottom=1.5mm]
\begin{tabularx}{\textwidth}{@{}>{\bfseries\color{AIEPNavy}}C{13mm} *{4}{>{\centering\arraybackslash}X}@{}}
& \textbf{M1} & \textbf{M2} & \textbf{M3} & \textbf{M4} \\
\hline
3° A & 07/09 & \cellcolor{GeoGoldSoft}14/09 & 21/09 & 28/09 \\
3° C & 09/09 & \cellcolor{GeoGoldSoft}23/09 & 30/09 & \cellcolor{GeoCyanSoft}\textbf{07/10} \\
\end{tabularx}
\end{tcolorbox}
{\small\color{AIEPMuted}El bloque de 3° A del 14/09 se mantiene. 3° C no tiene bloque el 16/09 y retoma el 23/09, conservando sus cuatro mentorías en orden.}

\vspace{3mm}
\begin{tcolorbox}[
  enhanced,
  colback=white,
  colframe=AIEPRed,
  boxrule=0.65pt,
  borderline west={2.2mm}{0pt}{AIEPRed},
  arc=1mm,
  left=4mm,right=4mm,top=2.4mm,bottom=2.4mm]
{\scriptsize\bfseries\color{AIEPRed}AJUSTE NECESARIO PARA EL CIERRE}\par
{\normalsize\bfseries\color{AIEPNavy}La instancia final común debe realizarse después del 7 de octubre.}\par
{\small\color{AIEPInk}La fecha anterior del 5 de octubre dejaría a 3° C sin su Mentoría 4 de presentación y ensayo. Situar el cierre después del último bloque asegura que ambos cursos completen la misma ruta formativa antes de presentar.}
\end{tcolorbox}

\vspace{2mm}
\begin{tcolorbox}[
  colback=AIEPSoft,
  colframe=AIEPLine,
  boxrule=0.45pt,
  arc=1mm,
  left=3mm,right=3mm,top=2mm,bottom=2mm]
{\scriptsize\bfseries\color{AIEPNavy}CONDICIONES DE EJECUCIÓN CONFIRMADAS}\par
{\small\color{AIEPInk}
Laboratorio por curso · proyector disponible · computadores con acceso a internet para estudiantes · AIEP aporta los demás materiales.}
\end{tcolorbox}

\end{document}
